% META: {"id": "TCOMPL-MF-EX-039", "chapitre": "TCOMPL-MODELES-FONCTION", "type_objet": "exercice", "capacites": ["TCOMPL-MF-C4"], "capacites_codes": ["C4"], "methodes": ["M4"], "parcours": 2, "competences": ["raisonner"], "duree_min": 15, "mode_creation": "ex_nihilo", "sources_inspiration": [], "fichier_tex": "chapitres/TCOMPL-MODELES-FONCTION/exercices/TCOMPL-MF-EX-039.tex", "corrige_tex": "chapitres/TCOMPL-MODELES-FONCTION/corriges/TCOMPL-MF-CO-039.tex", "status": "approved"}
% BEGIN-VERIFY
% from sympy import *
% x = symbols('x')
% f = ln(x)
% fp = diff(f, x)
% # tangente en x=1: f(1)=0, f'(1)=1, y=x-1
% # ln(x) - (x-1) = ln(x) - x + 1 <= 0 (concavite)
% g = ln(x) - x + 1
% g2 = diff(g, x, 2)
% assert g2 == -1/x**2  # < 0, concave
% assert g.subs(x, 1) == 0
% assert g.subs(x, 2) == ln(2) - 1 < 0
% assert simplify(g.subs(x, Rational(1,2)) - (-ln(2) + Rational(1,2))) == 0
% END-VERIFY
\begin{exercice}{TCOMPL-MF-EX-039}{2}{15}
Soit $f(x) = \ln x$ sur $\left]0, +\infty\right[$.
\begin{enumerate}
  \item Determiner la tangente a la courbe au point d'abscisse $1$.
  \item Montrer que la courbe est au-dessous de cette tangente.
  \item En deduire une inegalite connue.
\end{enumerate}

\end{exercice}
